\section{Problem Formulation}
\label{sec:formulation}

\subsection{Sequential Synthetic Dataset Construction}

Let $D_{\mathrm{train}}=\{(z_j,y_j)\}_{j=1}^{n}$ denote the real TRAIN split,
where $z_j$ contains the non-target fields and $y_j$ is the target. Let $S$
denote the dataset specification: field names and types, the target definition,
categorical domains, missing-value policy, and any declared logical or domain
constraints. For a requested class or generation condition $c$, a generator
$G$ proposes a candidate record
\begin{equation}
  x_t \sim G\!\left(\,\cdot \mid S,c,C_{\mathrm{train}}\right),
  \label{eq:generation}
\end{equation}
where $C_{\mathrm{train}}$ is context derived only from $D_{\mathrm{train}}$.
The real TEST split is excluded from $C_{\mathrm{train}}$, validator fitting,
and controller design.

The accepted synthetic dataset evolves with the decision sequence. Writing
$D_{\mathrm{syn}}^{(t)}$ for the collection before evaluating $x_t$, an
acceptance produces
\begin{equation}
  D_{\mathrm{syn}}^{(t+1)}
  =D_{\mathrm{syn}}^{(t)}\cup\{x_t\},
  \label{eq:accepted-update}
\end{equation}
whereas rejection leaves the collection unchanged. Repair and regeneration
create a new candidate version or a new candidate, respectively, and therefore
require another validation pass. This sequential view matters because a
candidate's dataset-level value depends on records accepted earlier.

\subsection{Quality, Risk, Uncertainty, and Contribution}

We distinguish four objects. Record-level quality is a vector
\begin{equation}
  Q_R(x_t)=
  \bigl(q_{\mathrm{schema}},q_{\mathrm{constraint}},
        q_{\mathrm{logical}},q_{\mathrm{semantic}},q_{\mathrm{support}}\bigr),
  \label{eq:record-quality}
\end{equation}
whose components need not be commensurate. A binary validity decision, when a
protocol requires one, is a policy over relevant components rather than a
definition of quality itself.

Record-level risk $R_R(x_t)$ concerns the expected consequence of accepting a
candidate under a stated use and loss. The same defect can therefore carry
different risk in exploratory augmentation and a high-stakes deployment.
Evaluator uncertainty $U_R(x_t)$ describes confidence, disagreement, missing
evidence, or limited applicability of the evaluators. It is not a synonym for
risk: an evaluator may be certain about a low-consequence defect or uncertain
about a high-consequence one.

Finally, dataset marginal contribution is conditional on the current accepted
set:
\begin{equation}
  \Delta_D\!\left(x_t\mid D_{\mathrm{syn}}^{(t)}\right)
  =\bigl(\Delta_{\mathrm{novelty}},
         \Delta_{\mathrm{coverage}},
         \Delta_{\mathrm{fidelity}},
         \Delta_{\mathrm{utility}}\bigr).
  \label{eq:dataset-contribution}
\end{equation}
These components may trade off and are not collapsed into one weighted score
in this draft. Individual validity does not imply positive marginal
contribution across all dataset-quality dimensions. An exact copy can be
structurally and semantically valid while adding no new support. Conversely,
novelty cannot excuse an invalid record.

\subsection{Scoped Heterogeneous Evidence}

For candidate $x$, let
\begin{equation}
  \mathcal{E}(x)=\{e_1(x),\ldots,e_m(x)\}
  \label{eq:evidence-set}
\end{equation}
be a set of evidence items. Conceptually, each item is a tuple
\begin{equation}
  e_i=(s_i,d_i,v_i,u_i,p_i),
  \label{eq:evidence-item}
\end{equation}
where $s_i\in\{\mathsf{FIELD},\mathsf{RECORD},\mathsf{SUBGROUP},
\mathsf{DATASET}\}$ is scope; $d_i$ is a quality or risk dimension; $v_i$ is
a score, verdict, or structured finding; $u_i$ describes confidence or
uncertainty; and $p_i$ records provenance. Evidence at different scopes is not
automatically comparable. A failed type check can identify a field-level
violation, while a coverage statistic describes the accepted dataset and
cannot be attributed to one field in the same way.

The current repository contract implements versioned evidence identifiers,
quality dimension, verdict, severity, optional score and confidence, affected
fields, provenance, latency, cost, and creation time. A dedicated scope field
is part of the proposed architecture and is not yet implemented. This
difference is recorded because the mathematical object describes the intended
M2+ interface rather than overstating the M1 contract.

\subsection{Decision Problem}

The intended controller applies a policy
\begin{equation}
  a_t=\pi\!\left(\mathcal{E}(x_t),
                  D_{\mathrm{syn}}^{(t)},S,\rho\right),
  \quad
  a_t\in\{\mathsf{ACCEPT},\mathsf{REPAIR},
           \mathsf{REGENERATE},\mathsf{REJECT}\},
  \label{eq:controller}
\end{equation}
where $\rho$ denotes declared risk constraints and operating preferences.
Revalidation is an internal transition required after repair; the repository
action contract reserves it explicitly. Repair defines a longer-range
extension of the action space and is not required for the first core
acceptance-decision comparison. The policy should retain abstention or
unknown states when evidence is missing and should not treat an uncalibrated
critic score as a correctness probability. Equation~\eqref{eq:controller}
defines the proposed decision problem. No learned fusion or risk-constrained
controller is implemented in the present milestone.
